\section{Representative applications}

The examples in this section are not all claimed as completed formalizations.
Their purpose is to show exactly how a classical argument is decomposed into
checked algebra, finite control, relational observations, and explicit
producer obligations.

\subsection{Polynomial identities and triangle formulas}

The simplest successful case is already represented by the checked
\path{IPRADerivations} library.  A geometric identity is translated into an
integer polynomial identity by representing signed quantities as pairs of
naturals.  The local theorem is generated by \path{ipra_ring}, universally
closed, and exported as an ordinary iPRA sentence.

In the present framework, such a theorem becomes an algebra node.  Further
finite reasoning---for example, substituting side-length expressions computed
by a coordinate trace, checking positivity hypotheses, or selecting among
finite combinatorial cases---can occur in the native graph.  The final graph is
exported again.  This demonstrates that direct derivations and clause programs
are complementary rather than competing designs.

\subsection{Fixed-dimensional Cayley--Hamilton}

Fix a dimension \(d\).  The theorem states that a \(d\times d\) matrix over a
commutative ring satisfies its characteristic polynomial.  Over integer-coded
entries, a finitary iPRA route is:

\begin{enumerate}
\item code matrices and polynomial matrices as fixed finite tuples;
\item define determinant by the finite Leibniz sum;
\item prove the adjugate identity
\[
   (tI-A)\operatorname{adj}(tI-A)=\det(tI-A)I
\]
as finitely many polynomial identities;
\item compare coefficients or substitute the matrix \(A\) in the polynomial
identity, using only finite sums and products;
\item discharge each entry equality by the ring backend;
\item bundle the \(d^2\) equations into the matrix conclusion.
\end{enumerate}

For fixed \(d\), all permutation lists and expression sizes are constants, so a
host generator can produce concrete \(\PTm\) expressions and
\path{ipra_ring} calls.  The generated theorem remains symbolic in the matrix
entries.  A graph node packages the final matrix equation.

For variable \(d\), additional generic libraries are needed, but the recursion
is still bounded by explicit finite permutation and index ranges.  The main
new obligation is proof engineering, not a foundational obstacle.

\subsection{A certified root-cluster computation}

Let \(p\in\Q[z]\) and let a query request root information to dyadic scale
\(2^{-q}\) in a bounded rectangle.  A robust theorem should not return an
ordering of all complex roots.  Instead, a producer supplies a finite family of
rational rectangles or disks with:

\begin{itemize}
\item a certified root count for each region;
\item pairwise disjointness;
\item total count equal to the degree;
\item diameter at most \(2^{-q}\);
\item refinement edges to the previous level.
\end{itemize}

The verifier checks all rational geometry and the root-count certificates.  The
particular root-count method may be a Sturm/subresultant argument for real
roots, an argument-principle winding computation with separation bounds, or a
Rouch\'e certificate.  The finite cloud contains all verified decompositions.
Backward viability chooses a compatible finite refinement path if needed.

Stable observers include:

\begin{itemize}
\item the multiset of cluster multiplicities;
\item symmetric polynomials of cluster centers with explicit error;
\item the number of roots in a fixed separated region;
\item an individual root approximation only for a singleton cluster separated
from all others by a positive gap.
\end{itemize}

At a multiple root, the theorem still returns a small cluster with multiplicity;
it does not fail by trying to name indistinguishable branches.

\subsection{Finite spectral information}

For a rational or represented-real matrix, exact eigenvectors are unstable near
collisions.  A stable contract requests:

\begin{enumerate}
\item a characteristic polynomial approximation or exact polynomial;
\item certified spectral clusters;
\item multiplicity or invariant-factor data;
\item when a gap \(\delta\) separates a cluster, an approximation to its
spectral projector or invariant subspace.
\end{enumerate}

Matrix arithmetic and characteristic polynomial identities use the algebra
backend.  Cluster construction uses relational nodes.  Projector formulas are
guarded by a resolvent or separation bound.  The final observer may be the
trace, determinant, rank, cluster dimension, or projector action on finitely
many vectors.  A basis is exported only when a verified normalization and pivot
scheme is stable under the supplied gap.

\subsection{Numerical integration with an explicit modulus}

Suppose a represented function \(f:[0,1]\to\R\) comes with:

\begin{itemize}
\item an evaluation contract returning rational approximations;
\item a uniform continuity modulus \(\omega\);
\item a magnitude bound \(M\).
\end{itemize}

To approximate the integral within \(2^{-q}\), choose a mesh exponent
\(m=\omega(q+c)\) and evaluate a rational Riemann sum.  The demand pass requests
function values at enough precision to control arithmetic error.  The local
proof splits the total error into oscillation on subintervals and evaluation
error.  All estimates are finite rational inequalities.  The sum is a bounded
fold, and polynomial rearrangements use the algebra backend.

The theorem is therefore constructive and uniform in the supplied modulus.  If
only classical continuity is given, the modulus producer remains an explicit
hypothesis or is replaced by a finite cover certificate on the requested
scale.

\subsection{Formal power series and analytic functions}

For a formal identity such as
\(\exp(x+y)=\exp(x)\exp(y)\) up to degree \(N\), the object is a finite
coefficient theorem.  Coefficients are rational with factorial denominators;
all divisions are by positive numerals.  Bounded convolution and binomial
identities prove the result.  No convergence theorem is needed.

For analytic evaluation, add a radius bound and a remainder estimate.  The
producer supplies or derives a tail bound.  The demand pass computes the
truncation order and coefficient precision.  The final theorem states a
rational error bound.  This separates algebraic coefficient identities from
analysis, a distinction that simplifies many Gamma, trigonometric, and
special-function packets.

\subsection{Finite topology and homology}

Given a finite simplicial complex, the chain groups and boundary matrices are
finite.  A Smith-normal-form certificate consists of matrices \(U,D,V\) with
\(U\) and \(V\) unimodular and
\(U\partial V=D\), plus divisibility conditions on the diagonal.  Verification
is finite integer arithmetic.  Homology ranks and torsion invariants follow by
an object-level algebraic finalizer.

For an ideal space, a separate producer supplies a finite complex and a theorem
relating it to the requested topological observation: a triangulation,
nerve equivalence, deformation retraction, or certified approximation.  The
framework fully handles the finite algebraic topology once that bridge is
present.  It does not claim that every topological space has a uniformly
computable triangulation.

\subsection{Finite categories and diagram chasing}

A finite diagram-chasing theorem is represented by object and morphism tables,
commutativity equations, and universal-property witness clauses.  Composition
is table lookup.  Naturality and uniqueness are bounded equations.  The graph
compiler mirrors the diagram: each commuting square or universal property is a
local clause node, and the terminal equality is obtained by cut and
substitution.

For example, uniqueness of an inverse or uniqueness in a pullback can be
proved with a small symbolic equality chain.  Yoneda component calculations
need only the finite morphisms appearing in the formula, not a search through
an entire category.

\subsection{Diophantine enumeration given a height bound}

Let \(F(\vec x)=0\) be a Diophantine equation and suppose a producer supplies a
bound \(H\) such that every solution in the desired domain has each coordinate
at most \(H\).  The remaining theorem is uniform:

\begin{enumerate}
\item enumerate the finite box \([0,H]^n\);
\item evaluate \(F\) by PR arithmetic;
\item retain exactly the zeroes;
\item prove soundness and completeness by bounded trace invariants;
\item compute any finite statistic or output the complete list.
\end{enumerate}

If the producer instead supplies a complete finite list, use the second half of
Theorem~\ref{thm:height-list}.  If only a cardinality bound is known, no finite
box follows.  This locates the exact gap in general Siegel or Faltings
formalizations.

\subsection{Prime-distribution statements}

A finite quantitative prime theorem at parameter \(x\) is often entirely
checkable: primality below \(x\), residue classes, weighted sums, and rational
inequalities are bounded computations.  A general asymptotic theorem needs a
function producing thresholds and constants for requested error.

The framework suggests a two-part statement:

\begin{enumerate}
\item \emph{certificate theorem}: given explicit constants, zero-free regions,
exceptional-character data, and thresholds satisfying bounded checks, the
claimed estimate holds for the finite range or all larger values by a proved
analytic recurrence;
\item \emph{producer theorem}: the required constants exist and are generated
uniformly from the parameters.
\end{enumerate}

The first part may be approachable in iPRA after a substantial finite analytic
library.  The second may inherit ineffectivity from the classical proof.  The
report does not merge the two.

\subsection{A general application template}

For a new headline theorem, prepare the following table before writing Lean.

\begin{longtable}{L{0.22\textwidth}L{0.69\textwidth}}
\toprule
Field & Required content \\
\midrule
Input presentation & Numerical codes, validity, and external adequacy.\\
Output query & Exact finite observation requested by a precision/resource term.\\
Producer data & Algorithms or explicit hypotheses supplying moduli, gaps, covers, heights, thresholds, or finite certificates.\\
Verifier & Bounded Boolean checks and their local correctness derivations.\\
Demand & Monotone backward transformer and resource envelope.\\
Control & DAG, bounded trace, descent, search, or relational viability.\\
Algebra & Direct \(\Prf\;\PRAAx\) lemmas, preferably from shared libraries.\\
Observer & Exact or approximate branch-independent output; gap if an exact label is claimed.\\
Final formula & Concrete \(\NFm\) or \(\mathrm{Fm}\;\PTm\) expression.\\
Endpoint & Terminal \(\KPrf\), \(\NPrf\), and/or \(\Prf\;\PRAAx\) type.\\
Status & Human proof, executable tests, Lean outline, and kernel-check status recorded separately.\\
\bottomrule
\end{longtable}
